\selectlanguage{spanish}

\renewcommand{\articuloidiomaxml}{es}
\renewcommand{\articulotitulo}{Las deudas pendientes de la digitalización en América Latina}
\renewcommand{\articuloautores}{Andrea Esther Liberatore}
\renewcommand{\articulodoi}{10.1787/7a127615-es}
\renewcommand{\articulotipo}{EDITORIAL}
\renewcommand{\articulorevista}{Agenda Digital Latinoamericana}
\renewcommand{\articuloissne}{1853-0028}
\renewcommand{\articulovolumen}{01}
\renewcommand{\articulonumero}{01}
\renewcommand{\articulofecha}{ 2026}
\renewcommand{\articulopagina}{1}

\setcounter{page}{1}
\thispagestyle{firstpage}

% --- TÍTULO A ANCHO COMPLETO ---
\noindent\begin{minipage}[t]{\dimexpr\textwidth+\marginparsep+\marginparwidth\relax}%
{\sffamily\bfseries\Large\articulotitulo\par}%
\end{minipage}\par
\vspace{3mm}%
\renewcommand{\articulotitulotrans}{The Pending Challenges of Digitalization in Latin America}
\noindent\begin{minipage}[t]{\dimexpr\textwidth+\marginparsep+\marginparwidth\relax}%
{\sffamily\itshape\normalsize\articulotitulotrans\par}%
\end{minipage}\par
\vspace{4mm}%

% --- BLOQUE LATERAL DE METADATOS ---
\begin{textblock*}{68mm}(\gbbloquex,92mm)
{\setlength{\leftskip}{0pt}\setlength{\parindent}{0pt}%
\noindent\begin{minipage}[t]{68mm}
\sffamily\small\setlength{\parskip}{1pt}
{\bfseries\color{azulrevista}Andrea Esther Liberatore}\par
{\scriptsize\texttt{\href{https://orcid.org/0009-0003-9648-5013}{https://orcid.org/0009-0003-9648-5013}}}\par
Departamento de Ciencias de Datos y Tecnología Social, Instituto de Estudios sobre Conectividad y Desarrollo Latinoamericano, Ciudad Autónoma de Buenos Aires, Argentina\par
\vspace{6pt}
{\bfseries\color{azulrevista!70}Derechos de autor\'ia:}\par
\textcopyright\ 2026 Liberatore, A. Publicado por Agenda Digital Latinoamericana. Se permite el uso, distribución y reproducción sin fines comerciales, siempre que se cite la obra original y las obras derivadas se distribuyan bajo la misma licencia. (\href{https://creativecommons.org/licenses/by-nc-sa/4.0/}{https://creativecommons.org/licenses/by-nc-sa/4.0/})\par\vspace{6pt}
{\bfseries\color{azulrevista!70}Publicación:} 2026, vol.~01, n\'um.~01, pp.~1--7\par\vspace{6pt}
{\bfseries\color{azulrevista!70}URL:} {\scriptsize\texttt{\href{https://www.clacso.org/ciencia-tecnologia-y-sociedad-en-america-latina/}{https://www.clacso.org/ciencia-tecnologia-y-sociedad-en-america-latina/}}}\par\vspace{6pt}
{\bfseries\color{azulrevista!70}DOI:} {\scriptsize\texttt{\href{https://doi.org/10.1787/7a127615-es}{10.1787/7a127615-es}}}\par\vspace{6pt}
{\bfseries\color{azulrevista!70}Licencia:} CC BY-NC-SA 4.0\par
{\scriptsize\texttt{\href{https://creativecommons.org/licenses/by-nc-sa/4.0/}{https://creativecommons.org/licenses/by-nc-sa/4.0/}}}\par
\vspace{6pt}
{\bfseries\color{azulrevista!70}Fechas}\par\vspace{2pt}
Recibido: 01/05/2026\par
Aceptado: 26/05/2026\par
Online: 26/06/2026\par
\vspace{6pt}
\vspace{6pt}
{\bfseries\color{azulrevista!70}Compartir:}\par\vspace{3pt}
{\large
\href{https://bsky.app/intent/compose?text=https%3A%2F%2Fwww.clacso.org%2Fciencia-tecnologia-y-sociedad-en-america-latina%2F}{\includesvg[height=0.9em]{/home/andrea/.gbpublisher/fonts/bluesky}}~
\href{mailto:?subject=Las%20deudas%20pendientes%20de%20la%20digitalizaci%C3%B3n%20en%20Am%C3%A9rica%20Latina&body=https%3A%2F%2Fwww.clacso.org%2Fciencia-tecnologia-y-sociedad-en-america-latina%2F}{\includesvg[height=0.9em]{/home/andrea/.gbpublisher/fonts/envelope}}~
\href{https://www.facebook.com/sharer/sharer.php?u=https%3A%2F%2Fwww.clacso.org%2Fciencia-tecnologia-y-sociedad-en-america-latina%2F}{\includesvg[height=0.9em]{/home/andrea/.gbpublisher/fonts/facebook}}~
\href{https://www.linkedin.com/sharing/share-offsite/?url=https%3A%2F%2Fwww.clacso.org%2Fciencia-tecnologia-y-sociedad-en-america-latina%2F}{\includesvg[height=0.9em]{/home/andrea/.gbpublisher/fonts/linkedin}}~
\href{https://www.reddit.com/submit?url=https%3A%2F%2Fwww.clacso.org%2Fciencia-tecnologia-y-sociedad-en-america-latina%2F&title=Las%20deudas%20pendientes%20de%20la%20digitalizaci%C3%B3n%20en%20Am%C3%A9rica%20Latina}{\includesvg[height=0.9em]{/home/andrea/.gbpublisher/fonts/reddit}}~
\href{https://www.researchgate.net/search?q=Las%20deudas%20pendientes%20de%20la%20digitalizaci%C3%B3n%20en%20Am%C3%A9rica%20Latina}{\includesvg[height=0.9em]{/home/andrea/.gbpublisher/fonts/researchgate}}~
\href{https://t.me/share/url?url=https%3A%2F%2Fwww.clacso.org%2Fciencia-tecnologia-y-sociedad-en-america-latina%2F&text=Las%20deudas%20pendientes%20de%20la%20digitalizaci%C3%B3n%20en%20Am%C3%A9rica%20Latina}{\includesvg[height=0.9em]{/home/andrea/.gbpublisher/fonts/telegram}}~
\href{https://wa.me/?text=https%3A%2F%2Fwww.clacso.org%2Fciencia-tecnologia-y-sociedad-en-america-latina%2F}{\includesvg[height=0.9em]{/home/andrea/.gbpublisher/fonts/whatsapp}}~
\href{https://twitter.com/intent/tweet?url=https%3A%2F%2Fwww.clacso.org%2Fciencia-tecnologia-y-sociedad-en-america-latina%2F&text=Las%20deudas%20pendientes%20de%20la%20digitalizaci%C3%B3n%20en%20Am%C3%A9rica%20Latina}{\includesvg[height=0.9em]{/home/andrea/.gbpublisher/fonts/x-twitter}}
}\par
\end{minipage}%
}% FIN GRUPO leftskip
\end{textblock*}


\begin{tcolorbox}[colback=brown!8!white,colframe=brown!8!white,boxrule=0pt,arc=2pt,left=4pt,right=4pt,top=3pt,bottom=3pt,before skip=4pt,after skip=4pt]
\begin{otherlanguage}{spanish}
\sffamily\small \textbf{Resumen:} Este artículo editorial examina el impacto de la pandemia de COVID-19 como acelerador de la transformación tecnológica en América Latina, analizando cómo la migración virtual profundizó las asimetrías socioeconómicas preexistentes. El texto presenta un diagnóstico dicotómico: reconoce la expansión de la conectividad pospandemia, pero denuncia la persistencia de una "brecha digital de segundo nivel", caracterizada por limitaciones de acceso material y falta de competencias digitales en sectores vulnerables. Se postula que la digitalización sin políticas públicas integrales actúa como un vector que reproduce desigualdades históricas de género, etnia y territorio. Finalmente, se introduce este número de la revista, cuyas contribuciones abordan la problemática bajo la premisa de que la transformación digital no debe aislarse de sus condiciones materiales.
\par\smallskip\noindent\textbf{Palabras clave:} pandemia de covid-19, transformación tecnológica, américa latina, migración virtual, asimetrías socioeconómicas, brecha digital de segundo nivel, competencias digitales, políticas públicas, desigualdades históricas
\end{otherlanguage}
\end{tcolorbox}

\begin{tcolorbox}[colback=brown!8!white,colframe=brown!8!white,boxrule=0pt,arc=2pt,left=4pt,right=4pt,top=3pt,bottom=3pt,before skip=4pt,after skip=4pt]
\begin{otherlanguage}{english}
\sffamily\small \textbf{Abstract:} This editorial examines the impact of the COVID-19 pandemic as an accelerator of technological transformation in Latin America, analyzing how virtual migration deepened pre-existing socioeconomic asymmetries. The text presents a dichotomous diagnosis: it recognizes post-pandemic connectivity expansion while denouncing the persistence of a "second-level digital divide," characterized by material access limitations and a lack of digital skills in vulnerable sectors. It argues that digitalization without comprehensive, intersectional public policies acts as a vector that reproduces historical inequalities of gender, ethnicity, and territory. Finally, this journal issue is introduced, featuring contributions that address the issue under the premise that digital transformation cannot be isolated from its material conditions.
\par\smallskip\noindent\textbf{Keywords:} covid-19 pandemic, technological transformation, latin america, virtual migration, socioeconomic asymmetries, connectivity expansion, second-level digital divide, digital skills, public policies, historical inequalities, digital transformation
\end{otherlanguage}
\end{tcolorbox}

\bigskip
\noindent\rule{\linewidth}{0.4pt}
\bigskip

La pandemia de COVID-19 funcionó como un acelerador involuntario de procesos que llevaban años en gestación. En cuestión de semanas, sistemas educativos enteros migraron a plataformas virtuales, trámites administrativos que exigían presencialidad se trasladaron a portales web y la telemedicina dejó de ser una promesa para convertirse en una necesidad imperiosa. Sin embargo, esa aceleración no hizo más que volver más visibles las fracturas que ya existían en el tejido social de la región, según la CEPAL \autocite{2566-CEPAL2021}.

América Latina llegó a la pandemia con una infraestructura digital profundamente desigual. Mientras las capitales y las grandes áreas metropolitanas disponían de conexiones de banda ancha razonablemente estables, vastas zonas rurales e incluso periferias urbanas carecían de conectividad básica. La brecha no era solo de acceso físico: se manifestaba también en las competencias digitales de la población, en la disponibilidad de dispositivos adecuados y en la capacidad institucional para sostener servicios en línea de manera continua y segura.

Tres años después del momento más agudo de la crisis sanitaria, el balance resulta ambivalente. Por un lado, se registró una expansión significativa de la conectividad: varios países de la región superaron tasas de penetración de Internet que parecían lejanas a comienzos de 2020. Por otro, las desigualdades de uso —lo que la literatura especializada denomina “brecha digital de segundo nivel”— no solo persistieron sino que, en algunos casos, se profundizaron. Los sectores de menores ingresos accedieron mayoritariamente a través de teléfonos móviles con planes de datos limitados, lo que restringía severamente sus posibilidades de participar en actividades educativas, laborales o de ejercicio ciudadano que demandan conexiones estables y pantallas de mayor tamaño.\footnote{Según estimaciones de la CEPAL, en 2022 el costo del servicio de
banda ancha para la población del primer quintil de ingresos
representaba en promedio el 14\% de su ingreso mensual en la región, una
proporción que vuelve inviable cualquier pretensión de universalidad
efectiva.}

El problema, conviene insistir, no es exclusivamente tecnológico. La digitalización que no va acompañada de políticas públicas integrales —que atiendan simultáneamente la conectividad, la formación, el diseño institucional y la producción de contenidos pertinentes— corre el riesgo de reproducir y amplificar las desigualdades preexistentes. En una región donde la matriz de la desigualdad opera simultáneamente por ingreso, género, etnia, territorio y edad, suponer que la mera provisión de infraestructura resuelve el acceso es una simplificación que ningún responsable de política pública debería permitirse.

Este número de la revista reúne contribuciones que abordan distintas aristas de esta problemática. Los trabajos que aquí se presentan comparten una premisa: la transformación digital de nuestras sociedades no puede analizarse al margen de las condiciones materiales y simbólicas en las que se despliega. Esperamos que su lectura contribuya a enriquecer un debate que, lejos de cerrarse, se vuelve más urgente con cada innovación tecnológica que promete —una vez más— resolver por sí sola las deudas históricas de la región.

